\documentclass[11pt]{article}
\usepackage[margin=1in]{geometry}
\usepackage{amsmath,amssymb,amsthm,booktabs,hyperref}
\hypersetup{colorlinks=true,linkcolor=blue,urlcolor=blue}
\newtheorem{theorem}{Theorem}
\newtheorem{corollary}[theorem]{Corollary}
\newtheorem{lemma}[theorem]{Lemma}
\newcommand{\PiR}{\Pi^{\mathrm R}}
\newcommand{\supp}{\operatorname{supp}}
\newcommand{\Lip}{\mathcal L}
\title{Right-aligned packets and a qualitative current criterion for Collatz}
\author{Research continuation of the supplied Collatz programme handoff}
\date{17 October 2026}
\begin{document}
\maketitle
\begin{abstract}
At depth one, grouping exact source lifts from the right moves every ungrouped
endpoint into the fixed set $\{5,7,11,13\}$. This gives an error bound valid at
every time horizon. A maximum principle then proves that finite current
certificates for these revised packets at every integer cutoff imply Collatz,
with no uniform bound on certificate size or positive cost. More generally,
finite currents with everywhere nonpositive residual suffice. This is an
exact reformulation of the remaining problem. A further sampling argument
reduces the universal target to five explicit dyadic rays. A disjoint annular
recursion describes the added measure when a cutoff is multiplied by 64.
We also characterize finite-current existence by orbit components and prove
that uniformly bounded certificates for individual packets cannot suffice;
no universal current construction, Collatz proof, or novelty claim is supplied.
\end{abstract}

\section{Definitions and the change in grouping}
Let $S(x)=(3x+1)/2^{v_2(3x+1)}$ on positive odd integers, with $S(1)=1$.
Let $\tau(x)$ be the first hitting time of $1$, possibly infinite, and
$h_L(x)=\min(L,\tau(x))$. Measures below are finite signed integer measures on
odd states greater than $1$. Pushforward $P$ discards mass arriving at $1$.
Write $H_L(\mu)=\langle\mu,h_L\rangle$.

For an integer $N\ge8$, set
\[
b=\lceil\log_2N\rceil,\qquad
n_0=\left\lfloor\frac{N-1}{4}\right\rfloor,\quad
n_1=\left\lfloor\frac{3N+1}{2}\right\rfloor,\quad
n_2=\left\lfloor\frac{3N+1}{4}\right\rfloor,
\]
and retain the handoff's source measure
\[
\nu_N(x)=\mathbf1_{x\le n_1}(3\mathbf1_{x\equiv2\ (3)}-1)
 +\mathbf1_{x\le n_2}(3\mathbf1_{x\equiv1\ (3)}-1).
\]
Its signed total has absolute value at most $8$. For
$Q_T(u)=\sum_{1\le x\le u,\ x\text{ odd}}h_T(x)$ and the corresponding
residue sums $R_{r,T}$, the discrepancy is
\[
\mathcal E_T(N)=R_{5,T}(n_1)-\tfrac13Q_T(n_1)
 +R_{1,T}(n_2)-\tfrac13Q_T(n_2),\qquad
3\mathcal E_T(N)=H_T(\nu_N).
\]

For a single exact valuation $a$, put $M_a=2^{a+1}$. The source and endpoint
lifts are $x_t=q_a+M_at$ and $y_t=e_a+6t$, where
\[
(q_a,e_a,t_0)=
\begin{cases}
((2^a-1)/3,\ 1,\ 1),&a\text{ even},\\
((5\cdot2^a-1)/3,\ 5,\ 0),&a\text{ odd}.
\end{cases}                                                    \tag{1}
\]
The start $t_0$ excludes source $1$ and endpoints already absorbed at $1$.
In component $i\in\{1,2\}$ retain $t_0\le t\le\ell_{a,i}$, where
$\ell_{a,i}=\lfloor(n_i-q_a)/M_a\rfloor$. Set
\[
s_{a,i}=\max(0,\ell_{a,i}-t_0+1),\qquad r_{a,i}=s_{a,i}\bmod3.
\]
Leave the first $r_{a,i}$ lifts ungrouped and partition all remaining lifts
into consecutive triples. Thus the final triple ends at the upper cutoff.
The packet sum over $1\le a\le b$ and both components is denoted $\PiR_N$.
This is a new grouping, distinct from the handoff's left-aligned packets.

\begin{theorem}[Fixed-core remainder]
For every integer $N\ge8$,
\[
P\nu_N=\PiR_N+\eta_N,\qquad
\supp\eta_N\subseteq\{5,7,11,13\},\qquad
\|\eta_N\|_1\le6b.                                           \tag{2}
\]
Every raw packet has the form
$c_0\delta_y+c_1\delta_{y+6}+c_2\delta_{y+12}$ with coefficients a
permutation of $(2,-1,-1)$. In particular, $\langle1,\PiR_N\rangle=0$.
Moreover, for every integer $T\ge1$,
\[
\left|3\mathcal E_T(N)-H_{T-1}(\PiR_N)\right|\le8+30b.       \tag{3}
\]
\end{theorem}
\begin{proof}
Formula (1) solves $3x+1=2^ay$ with odd $y$, using its least positive source
representative modulo $2^{a+1}$. As $M_a$ is a unit modulo $3$, three
consecutive sources visit all three residues and their coefficients sum to
zero. At most two left lifts remain in each component. They have endpoints
$5,11$ for odd $a$ and $7,13$ for even $a$, and their absolute coefficients
sum to at most $3$ per component.

There is no unaccounted large-valuation tail: $x\le n_1<2^{b+1}$ and
$a(x)>b$ imply
\[
0<S(x)=\frac{3x+1}{2^{a(x)}}<3.
\]
The endpoint is odd, so it equals $1$ and is discarded by $P$. Summing the
leftovers proves (2). The four core hitting times are
\[
\tau(5)=1,\quad\tau(7)=5,\quad\tau(11)=4,\quad\tau(13)=2.
\]
Hence $|H_L(\eta_N)|\le30b$ for every $L$. Finally,
\[
H_T(\nu_N)=\langle1,\nu_N\rangle+H_{T-1}(P\nu_N)
=\langle1,\nu_N\rangle+H_{T-1}(\PiR_N)+H_{T-1}(\eta_N),
\]
which proves (3).
\end{proof}

\section{The qualitative criterion: no uniform current cost is needed}
For functions on nonnegative integers define
\[
(\Lip f)(N)=f(n_0)+\tfrac13f(n_1)+\tfrac13f(n_2).
\]
\begin{lemma}[A strict comparison function]
For every integer $N\ge16$,
\[
\Lip(N^{3/2})<\frac{505}{512}N^{3/2}.                         \tag{4}
\]
\end{lemma}
\begin{proof}
The three ratios are at most $1/4$, $49/32$, and $49/64$, respectively.
Consequently
\[
\frac{\Lip(N^{3/2})}{N^{3/2}}
\le\frac18+\frac{343}{384\sqrt2}+\frac{343}{1536}
<\frac18+\frac{245}{384}+\frac{343}{1536}
=\frac{505}{512}<1,
\]
using $1/\sqrt2<5/7$.
\end{proof}

\begin{theorem}[Equivalent qualitative targets]
The following four statements are equivalent.
\begin{enumerate}
\item Every positive integer reaches $1$ under Collatz.
\item For every integer $N\ge8$, $\sup_{L\ge0}H_L(\PiR_N)<\infty$.
      The bound may depend arbitrarily on $N$.
\item For every integer $N\ge8$, there is a finite signed current $J_N$
      such that $\PiR_N=(I-P)J_N$.
\item For every integer $N\ge8$, there are finite signed measures $J_N,\zeta_N$
      such that $\PiR_N=(I-P)J_N+\zeta_N$ and $\zeta_N\le0$ pointwise.
\end{enumerate}
Neither a uniform cost bound nor an inverse-basin theorem is required.
\end{theorem}
\begin{proof}
Statement 1 implies 3: a finite packet has finitely many sources, so under
Collatz all are absorbed by some finite time $M$. Then
$J_N=\sum_{j=0}^{M-1}P^j\PiR_N$ is finite and telescopes to $\PiR_N$.
Statement 3 implies 4 by taking $\zeta_N=0$. For a finite current,
\[
H_L((I-P)J)=\langle1,J-P^LJ\rangle
=\sum_yJ(y)\mathbf1_{\tau(y)\le L}\le\|J^+\|_1.
\]
Since $h_L\ge0$, a nonpositive residual contributes at most zero. Thus 4
implies 2 without any bound uniform in $N$.

It remains to prove 2 implies 1. Define the exceptional indicator
$e(x)=\mathbf1_{\tau(x)=\infty}$, with $e(1)=0$, and let
\[
B(u)=\sum_{1\le x\le u,\ x\text{ odd}}e(x),\qquad
B_r(u)=\sum_{1\le x\le u,\ x\equiv r\ (6)}e(x).
\]
The indicator is invariant under $S$. Finite support and (2) give
\[
\lim_{L\to\infty}\frac{H_L(\PiR_N)}L
=\langle e,\PiR_N\rangle
=\langle e,P\nu_N\rangle
=\langle e,\nu_N\rangle,                                  \tag{5}
\]
because $e$ vanishes on the fixed core. Statement 2 makes (5) nonpositive.

Partitioning sources by first valuation gives the exact exceptional-count
identity
\[
B(N)=B(n_0)+B_5(n_1)+B_1(n_2)
=\Lip B(N)+\tfrac13\langle e,\nu_N\rangle
\le\Lip B(N).                                               \tag{6}
\]
For completeness, valuation $1$ maps $4j+3$ bijectively to $6j+5$;
valuation $2$ maps $8j+1>1$ to $6j+1\ge7$; and valuation at least $3$
has sources $4z+1$ with $z$ odd and $S(4z+1)=S(z)$. The exceptional small
cases $1$ and $5$ both have $e=0$, so (6) has no additive constant.

The odd seeds $1,3,5,7,9,11,13,15$ all terminate, so $B(N)=0$ for $N<16$.
Also $B(N)\le N$, hence $B(N)/N^{3/2}\to0$. If $B$ were nonzero, this
ratio would attain a positive global maximum $M$ at some integer $N\ge16$.
Using (6) and (4) at that maximizer yields
\[
B(N)\le\Lip B(N)\le M\Lip(N^{3/2})
<M N^{3/2}=B(N),
\]
a contradiction. Thus all odd seeds terminate, and removing initial powers
of two gives Statement 1.
\end{proof}

\paragraph{A second check for exact currents.}
There is a short independent proof of $3\Rightarrow1$. If $r$ is the least
nonterminating odd seed, take $N=\lceil(2r-1)/3\rceil$. Then $n_1$ is $r$
or $r+1$, and $n_2<r$. Therefore
\[
\langle e,\nu_N\rangle=3\mathbf1_{r\equiv2\ (3)}-1\in\{-1,2\}.
\]
The core has zero $e$-pairing, whereas every finite divergence has zero
$e$-pairing by invariance. This is impossible. This short proof uses exact
currents; the maximum principle above is what also permits nonpositive residuals.

\section{An exact constant error on dyadic cutoffs}
For $N=2^b$, $b\ge3$, the core remainder depends only on $b\bmod6$:
\[
\begin{array}{c|rrr}
b\bmod6&\eta_N(5)&\eta_N(7)&\eta_N(13)\\\hline
0&2&-1&0\\
1&-1&1&-1\\
2&-1&0&-2\\
3&-1&-1&0\\
4&-1&-2&2\\
5&2&-3&1
\end{array}                                                   \tag{7}
\]
There is no mass at $11$. To prove periodicity, increase $b$ and $a$ by $6$.
The cutoffs multiply by $64$, while $q_{a+6}=64q_a+21$ and
$M_{a+6}=64M_a$. Since $n_i-q_a$ is odd, the subtraction of
$21/(64M_a)$ does not change its quotient's floor. Also the shifted sources
are $64x_t+21\equiv x_t\pmod3$, so their coefficients are unchanged.
Thus each old valuation has the same lift count and residual coefficients
at the shifted valuation.
The six new valuations $a=1,\ldots,6$ contribute zero in total: odd $a$ have
no leftovers, and even $a$ contribute the vectors $(-2,1),(1,-2),(1,1)$
on $(7,13)$, which sum to zero. Direct arithmetic at $b=3,\ldots,8$ gives (7).

Since $\langle1,\nu_{2^b}\rangle=-1$, substitution of the four core hitting
times into (7) gives the sharper transfer bound
\[
\boxed{\quad
H_{T-1}(\PiR_{2^b})-12\le3\mathcal E_T(2^b)
\le H_{T-1}(\PiR_{2^b})+1,\quad T\ge1.
\quad}                                                       \tag{8}
\]
The equivalence above quantifies over every integer cutoff. The dyadic
computation alone does not establish its hypothesis; the next section proves
a separate sampling theorem.

\section{A sparse universal target: five dyadic rays}
\begin{theorem}[Five rays]\label{thm:rays}
Let $(a_0,a_1,a_2,a_3,a_4)=(12,14,16,18,21)$ and
$\mathcal G=\{a_i2^b:0\le i\le4,\ b\ge3\}$.
Each of Statements 2--4 of Theorem 3, restricted to $N\in\mathcal G$,
is equivalent to Collatz.
\end{theorem}
\begin{proof}
All odd seeds below $96$ terminate, as checked by exact finite iteration
in the accompanying script. Equations (5)--(6) give
$B(N)\le\Lip B(N)$ at each allowed cutoff under the bounded-work hypothesis.
For $i+3=5d+j$, $0\le j\le4$, put
$U=a_j2^{b+d}$. Then $U\ge3N/2$ when $N=a_i2^b$, because
\[
\begin{array}{c|rrrrr}
a_i&12&14&16&18&21\\\hline
a_j2^d&18&21&24&28&32\\
3a_i/2&18&21&24&27&63/2
\end{array}
\]
Every $N$ is divisible by $8$. Monotonicity therefore gives
\[
B(N)\le B(N/4)+\tfrac13B(U)+\tfrac13B(U/2).                 \tag{9}
\]
Each argument on the right belongs to $\mathcal G$ or lies below $96$.
Define a comparison weight on the rays by
\[
F(a_i2^b)=2^{(5b+i)/3}.
\]
Use the same formula for the few smaller ray points on the right of (9).
The ray representation is unique. Since $5d+j-i=3$,
\[
F(U)=2F(N),\quad F(U/2)=2^{1-5/3}F(N),\quad
F(N/4)=2^{-10/3}F(N).
\]
Thus the comparison contracts by
\[
\lambda=2^{-10/3}+\tfrac23(1+2^{-5/3})
<\tfrac1{10}+\tfrac89=\tfrac{89}{90}<1,                     \tag{10}
\]
using $2^{10}>10^3$ and $2^5>3^3$.

If $B$ is nonzero, it is positive at some allowed cutoff because the grid
is unbounded. Also $B(N)/F(N)\to0$ along the grid, since
$B(N)\le N$ and $F(N)$ grows as a positive, ray-dependent constant times
$N^{5/3}$. The ratio therefore attains a positive global maximum $C$ at
an allowed cutoff. For arguments below $96$, $B=0\le CF$.
Applying (9)--(10) at the maximizer gives
$B(N)\le C\lambda F(N)<CF(N)=B(N)$, a contradiction.
Conversely, Collatz supplies exact finite currents for every finite measure.
The implications from currents to bounded work are unchanged.
\end{proof}
The comparison depends on the ray as well as the scale. Replacing it by a
single power of $N$ loses this five-state contraction. No infinite ray has
yet been certified, and no minimality claim for the number of rays is made.

\section{Disjoint annular increments under a factor of 64}
\begin{theorem}[Exact annular recursion]\label{thm:annular}
If $8\mid N$ and $Q=9N/4$, then
\[
\PiR_{64N}=\PiR_N+A_N,\qquad \eta_{64N}=\eta_N,
\]
where $\supp\PiR_N\subset(1,Q)$ and $\supp A_N\subset(Q,64Q)$.
For odd $y$ in $2^jQ<y<2^{j+1}Q$, the coefficient $A_N(y)$ depends only
on $j$ and $y\bmod18$:
\[
\begin{array}{c|rrrrrr}
j\backslash(y\bmod18)&1&5&7&11&13&17\\\hline
0&1&0&1&0&-2&0\\
1&2&1&-1&1&-1&-2\\
2&3&-1&-3&2&0&-1\\
3&1&-3&-2&3&1&0\\
4&-1&-2&-1&1&2&1\\
5&0&-1&0&-1&0&2
\end{array}
\]
All other coefficients are zero. Each band has signed mass zero.
\end{theorem}
\begin{proof}
The cutoffs at $64N$ are $64n_i(N)$. For each old valuation, shifting
$a$ to $a+6$ replaces $q_a,M_a$ by $64q_a+21,64M_a$.
Since $n_i$ is even and $q_a$ is odd, this preserves the last lift index:
subtracting $21/64$ cannot cross the nearest lower multiple of $M_a$.
The shifted source is $64x+21\equiv x\pmod3$, so its coefficient is also
preserved. Thus the old packets and leftovers occur unchanged among
valuations $7,\ldots,b+6$ at the new cutoff.

For the six new valuations, write $N=8k$. In each component the odd
valuations $1,3,5$ have lift counts divisible by $3$. The even valuations
$2,4,6$ leave two lifts per component; the combined coefficients on
$(7,13)$ are $(-2,1),(1,-2),(1,1)$, respectively. Their sum is zero.
This proves the remainder identity and identifies $A_N$ with the full
pushforward contribution of the new six valuations.

For a possible endpoint $y>1$ and a valuation $a\le6$, the source is
$x=(2^ay-1)/3$; it is integral exactly when $2^ay\equiv1\pmod3$.
Its component coefficients are
$3\mathbf1_{x\equiv2\ (3)}-1$ and
$3\mathbf1_{x\equiv1\ (3)}-1$.
The corresponding endpoint cutoffs are $2^{7-a}Q$ and $2^{6-a}Q$.
These cutoffs are even, so no odd endpoint lies on a band boundary.
Below $Q$, all three allowed valuations of either parity occur in each
component. Their source residues cycle through all classes modulo $3$,
and their coefficients sum to zero. Above $64Q$, none occurs.
Within band $j$, the first component includes $a\le6-j$, and the second
includes $a\le5-j$, subject to the integrality condition. Evaluation modulo
$9$ gives the displayed table. Since $Q=18k$, each band comprises a whole
number of periods, and every row of the table sums to zero.

Finally, any original endpoint is at most $(3n_1(N)+1)/2=Q+1/2$.
It is odd and $Q$ is even, so it lies below $Q$.
\end{proof}

\begin{corollary}[A concrete infinite construction target]
It suffices to construct finite currents with nonpositive residual for
$A_{a_i2^{r+6k}}$, for every $i=0,\ldots,4$, $r=3,\ldots,8$, and $k\ge0$.
\end{corollary}
\begin{proof}
The 30 base measures $\PiR_{a_i2^r}$ have cutoffs at most $5376$ and
endpoints below $12096$. Exact finite iteration verifies convergence of
every odd seed below this endpoint bound within 98 odd steps. Therefore
each base has the finite exact current $\sum_{j=0}^{97}P^j\PiR_{a_i2^r}$.
Along each of the 30 ladders, sum this base current and the finitely
many increment certificates to obtain one for each $\PiR_N$ in
Theorem \ref{thm:rays}. Nonpositive residuals remain nonpositive under
addition. The theorem then gives Collatz.
\end{proof}
The base existence claim uses finite verified absorption; the 30 full
current tables have not been stored. The infinitely many increments remain
uncertified.
The annular recursion is an identity of measures; it supplies no current
for $A_N$ and no decreasing termination parameter.

\section{What finite currents detect}
The following observation applies to any deterministic map with one
absorber, with pushforward killed there. Two states are in the same
\emph{orbit component} when their forward trajectories meet, allowing
unequal arrival times. A component is absorbing if its states reach the
absorber.

\begin{theorem}[Component criterion]
For a finite signed measure $\mu$, let $m_C$ be its total mass in an
orbit component $C$. A finite exact current for $\mu$ exists if and only if
$m_C=0$ in every nonabsorbing component. A finite current with nonpositive
residual exists if and only if $m_C\le0$ in each such component.
\end{theorem}
\begin{proof}
Every nonabsorbing component indicator is invariant, so its pairing with
a finite divergence vanishes. This proves necessity.
For sufficiency, in each component meeting $\supp\mu$, route the finitely
many sources to a common forward descendant, with arbitrary arrival
times. In the absorbing component route them to the absorber. The sum of
these finite path currents has residual $m_C\delta_{z_C}$ at each chosen
nonabsorbing descendant $z_C$, which proves both assertions.
\end{proof}

For one fixed measure, bounded work is weaker: beyond the largest finite
source hitting time,
\[
H_L(\mu)=L\sum_{C\ {\rm nonabsorbing}}m_C+
          \sum_{\tau(x)<\infty}\mu(x)\tau(x).
\]
It only requires the \emph{sum} of nonabsorbing masses to be nonpositive.
For example, on negative odd integers the cycles $-1\mapsto-1$ and
$-5\mapsto-7\mapsto-5$ give $H_L(\delta_{-1}-\delta_{-5})=0$,
but no such current exists. This example is outside the positive Collatz
domain; it does not claim a positive nonterminating orbit.

A complete semidecision search exposes progressively longer paths from all
sources, retaining historical intersections. In each finite undirected
component of exposed arrows, route source mass along a spanning tree to
one anchor (the absorber when present). Traversing an arrow backwards
uses a negative current coefficient. If all remaining component masses
are zero, or nonpositive in the one-sided version, this is a certificate.
The component proof supplies finite witness paths whenever a certificate
exists, so sufficiently deep exposure succeeds. A finite depth cap may
return unresolved. Universal termination remains unproved.

\begin{theorem}[No bounded per-packet catalogue]
For every integer $H\ge0$, put $K=2H+7$, $y=9\cdot2^K-1$, and
$N=4\cdot2^K+5$. The actual final valuation-one packet in the first
source component at this cutoff is
\[
\mu_y=2\delta_y-\delta_{y+6}-\delta_{y+12}.
\]
Any finite exact current for $\mu_y$, or one with nonpositive residual,
has more than $H$ nonzero rows, including when arbitrary auxiliary states
are allowed.
\end{theorem}
\begin{proof}
At valuation one the lifts are $x_t=3+4t$, $y_t=5+6t$.
For $t=3\cdot2^{K-1}-1$, the last source index is $t+2$.
Right alignment leaves two initial lifts, and the final coefficients are
$(2,-1,-1)$ because $t\equiv2\pmod3$.

The three starting points are perturbations by $9\cdot2^K$ of $-1,5,11$.
Their base trajectories are the fixed point $-1$ and
$5\to1$, $11\to17\to13\to5\to1$. Use the usual formula, fixing $1$,
to describe these base trajectories only. Their accumulated valuations
$A_j$ obey $A_j\le2H+2$ for $0\le j\le H$. Exact branch arithmetic gives
\[
S^j(z_0+9\cdot2^K)=z_j+9\cdot3^j2^{K-A_j}.
\]
At least five bits remain in the perturbation. The first orbit is therefore
$31\bmod32$, whereas the other two are in $\{1,5,11,13,17\}\bmod32$,
through time $H$. They cannot meet each other across these two groups,
even at unequal times, and the positive orbit cannot reach $1$.

If a current had at most $H$ nonzero arrows, its finite undirected component
containing $y$ would have to contain a negative source or the absorber:
otherwise summing the divergence identity there would give
$2=\sum_C\zeta\le0$. A path to such a vertex uses at most $H$ arrows.
In a deterministic graph, an undirected simple path is directed toward a
common meeting vertex, since an internal vertex cannot have two outgoing
arrows. Its endpoints therefore meet with each arrival time at most $H$,
contradicting the established separation.
\end{proof}
The obstruction concerns individual packets. It leaves open cancellation
between packets and recursively expanded certificates of unbounded size.

\section{What this changes, and what remains open}
The handoff's left-aligned remainder was paid for by a time horizon. That
created a source/time-domain issue and motivated a uniform positive current
cost target. The revised depth-one grouping places the remainder inside a
verified finite core. Theorems 1--3 consequently remove both requirements
from the logical route to Collatz.

The remaining target is precise: construct a finite current for $\PiR_N$,
possibly with an everywhere nonpositive residual, along the five rays in
Theorem \ref{thm:rays}. Alternatively, handle their annular increments in
Theorem \ref{thm:annular}; the 30 finite base cases have been discharged.
A proof of termination of that construction is essential.
Its running time, number of rows, and positive cost may be arbitrarily large
functions of $N$. A small unresolved \emph{positive} residual is insufficient
for this all-horizon criterion unless its own future contribution is bounded.
The packet work itself need not be nonpositive: direct finite iteration gives
$H_{39}(\PiR_{32})=57$. The target is a finite upper bound for each cutoff.

This is an equivalence with Collatz, so it does not establish that the
remaining construction is easy. The archive's currents refer to different
packet groupings and named finite cutoffs; they cannot be substituted for a
universal construction of $\PiR_N$.

\section{Audit and reproducibility}
The uploaded archive's SHA-256 matched the handoff:
\begin{center}\small\texttt{025c8fafe25940951fe408b32b1ec806259725a6b4f8855d9f2b022911fe5447}\end{center}
All 31 manifest entries matched. The Python stored-certificate verifier and
freshly compiled independent C++ verifier both passed for $N=4096$ and
$N=262144$, reconstructing packets, divergences, exit profiles, and the
reported finite bounds. Those are checks of the original grouping.

The companion \texttt{right\_aligned\_packets.py} independently constructs the
new packets and compares them with direct one-step pushforward for every
$8\le N\le2048$ and selected larger cutoffs. It also checks the dyadic core
formula through $b=1000$, the core hitting times, and the exact rational
comparison constant. Results are saved in \texttt{right\_aligned\_audit.json}.
These computations check the implementation and examples. The general
claims above rest on the written proofs and have not received human peer
review or formal proof-assistant verification.

The continuation scripts in
\texttt{experiments/current\_construction/} check the five-ray arithmetic,
all odd seeds below $65536$, the 30 finite bases, and 66 complete annular
identities through expanded cutoff $262144$. Independent finite-graph
checks cover 1728 map/measure cases in both certificate modes, and the
packet obstruction is checked at selected depths through 128.
The historical-intersection constructor also reproduces ten whole-packet
examples and stored currents at $N=2053$ (nonpositive residual) and
$N=65536$ (exact). At $2053$, some cancelling source components have not
been linked to $1$ in the exposed graph; their unknown common future is
therefore eliminated by the certificate. These remain finite examples.
The unified command is \texttt{python3 scripts/verify\_all.py --cpp}.
\end{document}
